\SourceLeaf{268}
\begin{list}{}{\setlength{\leftmargin}{1.5em}\setlength{\itemindent}{-1.5em}\setlength{\itemsep}{0pt}\setlength{\parsep}{0pt}}
\item \(e\) die Excentricität der Ellipse, durch deren Umdrehung um ihre kleine Achse die ellipsoidische Fläche erzeugt wird;
\item \(t\) und \(90^0-w\) die Länge und Breite eines unbestimmten Punkts auf dieser Fläche, mithin \(w\) den Winkel einer in diesem Punkte gegen die Fläche gezogenen Normale mit der kleinen Achse;
\item \(T\) und \(90^0-U\) die Länge und Breite des entsprechenden Punkts auf der Kugelfläche;
\item \(i\) die imaginäre Einheit \(\BareRoot{-1}\);
\item \(f\) die Charakteristik für eine willkürlich zu wählende Function.
\end{list}
\noindent Die Logarithmen sind immer die hyperbolischen.\par
\noindent Durch \(m\) wird das Vergrösserungsverhältniss bezeichnet werden, so verstanden, dass jedes Linearelement auf der ellipsoidischen Fläche sich zu dem entsprechenden Linearelement auf der Kugelfläche verhält wie 1 zu \(m\): dieses Verhältniss ist an jeder Stelle der einen und der andern Fläche ein bestimmtes, für verschiedene Stellen veränderlich.

Die einfachste Auflösung erhält man, indem man die willkürliche Function schlechthin ihrem Argumente gleich, oder
\[
f\upsilon=\upsilon
\]
setzt, und diese Übergangsart ist in der That auch die geeignetste, wenn die \textit{ganze} Oberfläche des Ellipsoids auf die Kugelfläche übertragen werden soll. Für die Anwendung auf geodätische Rechnungen, wo immer nur ein vergleichungsweise sehr kleiner Theil der Erdfläche in Betracht kommt, ist es aber, wie schon a. a. O. bemerkt ist, viel vortheilhafter, der Function noch einen constanten und zwar imaginären Theil beizufügen, oder
\[
f\upsilon=\upsilon-i\log k
\]
zu setzen. Es lassen sich dann der Halbmesser der Kugel und die Constante \(k\) so bestimmen, dass die das Vergrösserungsverhältniss ausdrückende Grösse \(m\), von deren geringer Ungleichheit innerhalb der Grenzen der dargestellten Fläche die Bequemlichkeit der Anwendung auf geodätische Rechnungen vornehmlich abhängt, für den mittlern Parallelkreis \(=1\), und bis zu einigen Graden Entfernung nach Norden und Süden kaum merklich von 1 verschieden wird; die Abweichung von dem Werthe 1 ist nemlich von der zweiten Ordnung in Beziehung
